% Methods

\subsection{Model~A: Open (Boundary-Driven) Simulation}
\label{sec:model-a}
In Model~A, edge agents observe real market data---specifically, one-minute OHLC bars
for the QQQ ETF (Invesco QQQ Trust, 2024 calendar year, $T = 23{,}790$ time steps).
At each step~$\tau$: (i)~the typical price
$P_\tau = (H_\tau + L_\tau + C_\tau)/3$ is computed; (ii)~each edge agent evaluates
Eq.~\eqref{eq:edge-state} using the pre-computed EMA and volatility from
Eqs.~\eqref{eq:ema}--\eqref{eq:ema-vol}; (iii)~all interior agents update
simultaneously via the CA rule. The grid is initialized to all-Hold ($s=1$).

For the heterogeneous (crash) model, a fraction~$f$ of interior agents are assigned
Rule~19286 (contrarian) uniformly at random, while the remaining $1-f$ agents use
Rule~19280 (majority). The minority agent positions are fixed for the duration of the
simulation (quenched disorder).

\paragraph{Simulation parameters.}
The default grid is $256 \times 256$ ($N = 65{,}536$ agents). Edge agents comprise
25\% of boundary cells ($f_e = 0.25$, yielding ${\sim}512$ edge agents), selected
with a fixed random seed for reproducibility. The EMA learning rate is
$\alpha = 1/60$ and the band width is $k = 1.0$.

\subsection{Model~B: Closed (Endogenous Price) Simulation}
\label{sec:model-b}
In Model~B, the price is generated endogenously from the grid's net order imbalance.
Let $N_{\text{buy},\tau}$ and $N_{\text{sell},\tau}$ denote the number of interior
agents in the Buy and Sell states at step~$\tau$, respectively. The net order
imbalance is:
\begin{equation}
  \varphi_\tau = \frac{N_{\text{buy},\tau} - N_{\text{sell},\tau}}{N_{\text{interior}}},
  \label{eq:imbalance}
\end{equation}
and the price updates as:
\begin{equation}
  P_{\tau+1} = P_\tau \cdot \exp\!\bigl(\lambda\,\varphi_\tau\bigr),
  \label{eq:price-update}
\end{equation}
where $\lambda > 0$ is the price impact parameter. Edge agents then evaluate
Eq.~\eqref{eq:edge-state} using the endogenous price $P_{\tau+1}$ with online
EMA updates, closing the feedback loop.

\subsection{Phase Sweep Protocol}
\label{sec:phase-sweep}
To map the crash transition, we sweep the minority fraction $f$ from $0.0$ to $0.50$.
For each value of~$f$, the full $T$-step price series drives a $128 \times 128$ grid,
and we record the sell fraction at every time step. Three summary statistics are
computed: mean sell fraction, peak sell fraction, and crash frequency (fraction of
steps where sell fraction exceeds~$0.50$). A fine-grained sweep with
$\Delta f = 0.01$ is performed around the transition region $f \in [0.08, 0.22]$
to resolve the critical threshold.
